% Standalone problem document, generated from the adjacent README.md.
% Compile with XeLaTeX. Mathematical content is maintained in Markdown.
\documentclass[11pt,a4paper]{article}
\usepackage[fontsize=10.5pt]{scrextend}
\usepackage[margin=22mm,headheight=15pt,headsep=9mm,footskip=11mm]{geometry}
\usepackage{fontspec}
\setmainfont{texgyrepagella}[Extension=.otf,UprightFont=*-regular,BoldFont=*-bold,ItalicFont=*-italic,BoldItalicFont=*-bolditalic]
\setsansfont{texgyreheros}[Extension=.otf,UprightFont=*-regular,BoldFont=*-bold,ItalicFont=*-italic,BoldItalicFont=*-bolditalic]
\setmonofont{texgyrecursor}[Extension=.otf,UprightFont=*-regular,BoldFont=*-bold,ItalicFont=*-italic,BoldItalicFont=*-bolditalic,Scale=MatchLowercase]
\usepackage{amsmath,amssymb}
\usepackage{unicode-math}
\setmathfont{texgyrepagella-math.otf}
\usepackage{xcolor,microtype,enumitem,needspace,fancyhdr}
\usepackage{bookmark}
\definecolor{ink}{HTML}{17344B}
\definecolor{accent}{HTML}{197D80}
\definecolor{muted}{HTML}{596773}
\hypersetup{colorlinks=true,linkcolor=accent,urlcolor=accent,pdftitle={PF-01 - Exact
positive semidefinite rank of subset-intersection
matrices},pdfauthor={Open Problems in Numerical Linear Algebra}}
\urlstyle{same}
\setlength{\parindent}{0pt}
\setlength{\parskip}{5pt plus 1pt minus 1pt}
\linespread{1.04}
\setlength{\emergencystretch}{3em}
\widowpenalty=10000
\clubpenalty=10000
\displaywidowpenalty=10000
\allowdisplaybreaks[1]
\setlist{nosep,leftmargin=1.5em}
\providecommand{\tightlist}{\setlength{\itemsep}{0pt}\setlength{\parskip}{0pt}}
\providecommand{\pandocbounded}[1]{#1}
\pagestyle{fancy}
\fancyhf{}
\fancyhead[L]{\sffamily\footnotesize\color{muted}OPEN PROBLEMS IN NUMERICAL LINEAR ALGEBRA}
\fancyhead[R]{\sffamily\bfseries\footnotesize\color{accent}PF-01}
\fancyfoot[L]{\parbox[b]{0.9\textwidth}{\sffamily\fontsize{8}{10}\selectfont\color{muted}Editorial ratings; a bounded search cannot certify that a problem remains unsolved.\\Literature check: 2026-09-10}}
\fancyfoot[R]{\sffamily\footnotesize\color{muted}\thepage}
\renewcommand{\headrulewidth}{0.3pt}
\renewcommand{\headrule}{\hbox to\headwidth{\color{accent}\leaders\hrule height \headrulewidth\hfill}}
\makeatletter
\renewcommand{\section}{\Needspace{5\baselineskip}\@startsection{section}{1}{0pt}{12pt plus 2pt minus 2pt}{4pt}{\sffamily\bfseries\large\color{ink}}}
\renewcommand{\subsection}{\Needspace{4\baselineskip}\@startsection{subsection}{2}{0pt}{9pt plus 1pt minus 1pt}{3pt}{\sffamily\bfseries\normalsize\color{ink}}}
\makeatother
\setcounter{secnumdepth}{0}
\begin{document}
{\sffamily\small\color{muted}Nonnegative and positive
factorizations\par}
\vspace{3pt}
{\raggedright\hyphenpenalty=10000\exhyphenpenalty=10000\sffamily\bfseries\fontsize{21}{24}\selectfont\color{ink}Exact
positive semidefinite rank of subset-intersection matrices\par}
\vspace{7pt}
{\sffamily\small\textbf{Difficulty:} challenging\quad\textbf{Importance:} interesting
to specialist\par}
{\sffamily\small\bfseries\color{accent}Status: Open\par}
\vspace{4pt}
{\color{accent}\hrule height 0.8pt}
\vspace{6pt}
\textbf{Rating rationale:} Challenging because the symmetric subset
structure has not yielded matching PSD-rank bounds; specialist
importance reflects this explicit family's role as a benchmark for
semidefinite factorization.

\textbf{Area:} structured positive semidefinite factorization

\subsection{Context and notation}\label{context-and-notation}

Write \(\mathbb S_+^k\) for the cone of real symmetric positive
semidefinite \(k\times k\) matrices. For an entrywise nonnegative matrix
\(M\in\mathbb R_+^{p\times q}\), its \textbf{real positive semidefinite
rank} is

\[
\operatorname{rank}_{\rm psd}(M)
=\min\{k\ge1:\ \exists A_1,\ldots,A_p,B_1,\ldots,B_q\in\mathbb S_+^k,\quad
M_{ij}=\operatorname{tr}(A_iB_j)\ \text{for all }i,j\}.
\]

This definition concerns a family of matrix factors indexed by the rows
and columns of \(M\).

\subsection{Problem statement}\label{problem-statement}

For each integer \(n\ge5\), let \(\mathcal I_n\) and \(\mathcal J_n\)
consist of the subsets of \(\{1,\ldots,n\}\) of cardinalities
\(\lfloor n/2\rfloor\) and \(\lceil n/2\rceil\), respectively. Define
the nonnegative matrix

\[
M^{(n)}\in\mathbb R_+^{\mathcal I_n\times\mathcal J_n},
\qquad M^{(n)}_{I,J}=|I\cap J|.
\]

\subsubsection{Question}\label{question}

Determine \(\operatorname{rank}_{\rm psd}(M^{(n)})\) exactly as a
function of \(n\), with real symmetric factors as defined above. The
family is one problem; its individual orders are not separate catalog
entries.

The source supplies the bounds

\[
\left\lceil\frac{\sqrt{1+8n}-1}{2}\right\rceil
\le \operatorname{rank}_{\rm psd}(M^{(n)})
\le 2\lceil\sqrt n\rceil.
\]

These matrices provide structured benchmarks for algorithms seeking
small positive semidefinite factorizations.

\subsection{References}\label{references}

Hamza Fawzi, João Gouveia, Pablo A. Parrilo, Richard Z. Robinson, and
Rekha R. Thomas, \href{https://arxiv.org/html/1407.4095}{\emph{Positive
semidefinite rank}}, Mathematical Programming \textbf{153} (2015),
133--177, §9.1, Problem 9.2. Arnaud Vandaele, François Glineur, and
Nicolas Gillis, \href{https://arxiv.org/pdf/1707.07953}{\emph{Algorithms
for Positive Semidefinite Factorization}}, Computational Optimization
and Applications \textbf{71} (2018), 193--219, §4.2, preprint
pp.~10--11.

\subsection{Status check --- 2026-09-10}\label{status-check-2026-09-10}

Rechecked \href{https://arxiv.org/html/1407.4095}{Fawzi et al., Problem
9.2} and
\href{https://arxiv.org/pdf/1707.07953}{Vandaele--Glineur--Gillis,
§4.2}, and searched by subset intersections, Johnson schemes and the
problem number. The sources support the displayed real-factor bounds but
not an exact formula. No later determination was located; the 2018
algorithm paper remains the latest explicit status source checked, so
the later evidence is limited to this search.
\end{document}
